\section{Introduction}
\label{sec:intro}

Humans come to know an unfamiliar physical world by acting in it. They push, drop, and collide objects, form an understanding of the principles behind what they see, and then carry that understanding back to vision: a glance at a new scene is enough to picture how it will unfold. We ask whether multimodal agents can learn the physical world through interaction in the same way. Improving on a practiced task is not enough to show this. An agent can improve over a few trials by searching a small action space, replaying a trajectory that worked, or exploiting a visual cue that happens to correlate with success. Learning in the sense we study means that what the agent gathers by interacting with a world becomes an understanding of that world, one that still holds when the agent looks at a new scene from it.

Current benchmarks leave this distinction partly hidden. Passive physical-reasoning benchmarks ask models to judge physical concepts or outcomes from fixed images, videos, or questions; they do not test what feedback from action changes. Interactive physics benchmarks such as PHYRE~\cite{phyre} and I-PHYRE~\cite{iphyre}, and their recent VLM-agent counterparts~\cite{deepphy,ipr1}, measure task success, but success on the acted-on instance can mix physical inference with search. Indeed, VLMs trained through interaction improve on their training task without acquiring physical intuitions that transfer to related tasks~\cite{schulzebuschoff2026}. Scientific-discovery benchmarks~\cite{physgym,newtonbench,discoverphysics} let agents probe worlds with hidden or altered laws, but their observations and outputs are numeric or symbolic. Visual benchmarks of hidden physics~\cite{comphy,inphyre} test transfer to new scenes, but supply the evidence instead of letting the agent collect it. What remains open is the full path from interaction to vision: whether an agent that has experimented in a physical world can foresee how a new scene of that world will look.

Our benchmark is built on one observation: interacting with a physical world and seeing it correctly are two sides of the same visual experiment. An experiment pairs a rendered scene, an intervention, and a time interval with the rendered outcome. During \emph{interaction}, the agent chooses experiments and the world shows their outcomes. During \emph{evaluation}, the roles reverse: the world poses a new experiment on an unseen scene, and the agent must understand the scene and predict the next frame of the world. Because the two phases share the same scenes, interventions, and renderer, the evaluation asks the agent to use what it learned in exactly the form in which it learned it.

Each world departs from canonical central-force physics through one of two equations: through a structural change to the force law, such as a different distance exponent or mass coupling, or through a hidden coefficient field that varies with distance or time. Prior knowledge of textbook physics is therefore not enough. The agent spends a budget of experiments from a public menu, sees each outcome as a rendered strobe frame, and at fixed points predicts the outcome of a host-posed checkpoint before seeing it. After interaction, all tools and feedback are frozen, and the agent answers held-out questions on new scenes from its interaction history alone.

Sibling worlds are what make this a test of learning. Worlds of the same equation share their scenes and experiments but follow different hidden laws, so the same scene unfolds differently in each. A correct prediction therefore requires knowledge of the particular world the agent explored, which neither textbook physics nor familiarity with the setting provides. Learning is then measured as what interaction changes, through evidence arms on identical questions. We organize the analysis around three questions. Does interaction improve how the agent sees new scenes (RQ1)? Is the improvement specific to the world's own law, so that another world's history fails, and does it differ between structural and field departures (RQ2)? Where does learning break: in choosing experiments, in forming a law from them, or in applying it to images (RQ3)? A yoked observer that sees the same history without having acted, and a law-in-text ceiling, answer the last question.

Our contributions are:
\begin{itemize}
  \item A benchmark that frames learning the physical world through interaction as a visual experiment seen from two sides: the agent experiments and observes during interaction, then foresees rendered outcomes of new experiments with interaction frozen.
  \item A law-contrastive evaluation built from sibling worlds that share their scenes but differ in their hidden law, so that a correct prediction requires knowledge of the explored world, together with an identifiability check that removes hidden laws the images cannot reveal.
  \item An evaluation protocol that measures learning as what interaction changes, through online learning curves and paired evidence arms (no history, own world, sibling world, yoked observer, own record, law in text), applied to four frontier multimodal agents.
  \item Findings on current agents. Interaction teaches the law of the world explored: it raises the strongest system from \V{GptSolNoneAcc}\% to \V{GptSolOwnAcc}\%, and another world's history pulls it to that world's answer. Yet interaction recovers less than half of what the stated law provides (\V{GptSolLawAcc}\% with the law in text), and the gap lies in forming the law from evidence, not in applying it to images. Having run the experiments gives no advantage over observing them.
\end{itemize}

\begin{figure*}[t]
  \centering
  % \includegraphics[width=\textwidth]{figures/fig1_hero.pdf}
  \fbox{\parbox[c][6cm][c]{0.95\textwidth}{\centering
    \textbf{Figure 1 placeholder (hero figure).}\\[0.5em]
    \emph{(a) World:} world family $\rightarrow$ equation (structural change or hidden field) $\rightarrow$ hidden-law variant $\rightarrow$ initial state, for central-force motion.\\
    \emph{(b) Interact:} the agent chooses experiments and sees rendered outcomes; checkpoint predictions trace a learning curve.\\
    \emph{(c) Foresee:} interaction is frozen; from a new initial state, the agent understands the scene and predicts the next frame of the world; the alternatives it is scored against follow other laws, and textbook physics is always wrong.
  }}
  \caption{\textbf{Learning a physical world through interaction.} Interaction and evaluation are the same visual experiment seen from two sides. During interaction, the agent poses experiments and the world shows their outcomes. During evaluation, the world poses new experiments on unseen scenes, and the agent must foresee their outcomes from what it learned.}
  \label{fig:hero}
\end{figure*}

Section~\ref{sec:related} positions the benchmark relative to prior work. Section~\ref{sec:benchmark} defines the visual experiment, the world hierarchy, and the episode. Section~\ref{sec:construction} describes the construction of worlds, tools, and renders. Section~\ref{sec:protocol} defines the metrics and controls. Section~\ref{sec:validation} gives validation checks, Section~\ref{sec:experiments} describes the experiments and analyses, and Section~\ref{sec:conclusion} discusses limitations and conclusions.
